% ==============================================================================
% SECCIÓN 7: VALIDACIÓN Y SIMULACIÓN COMPUTACIONAL DE RESULTADOS
% ==============================================================================

\section{Validación y Simulación Computacional de Resultados}
\label{sec:validacion_simulacion}

Con el firme propósito de garantizar la rigurosidad científica, la consistencia cinemática y la viabilidad mecánica del torno didáctico para madera, la totalidad de los veinte cálculos dinámicos formulados en la Sección~\ref{sec:calculos_dinamicos} fueron sometidos a un protocolo exhaustivo de validación computacional dual. Este marco metodológico comprende dos plataformas complementarias de simulación de alta fidelidad:
\begin{enumerate}[leftmargin=18pt, itemsep=4pt, parsep=2pt]
    \item \textbf{Simulación Multicuerpo 2D de Dinámica Restringida (Working Model 2005):} Modelado de cuerpos rígidos planos bajo restricciones cinemáticas de rotación pura alrededor de eje fijo baricéntrico, amortiguamiento viscoso interfacial y contacto mecánico. Permite comprobar de manera directa y visual las magnitudes escalares y vectoriales de velocidades en la periferia e interiores, aceleraciones intrínsecas (tangencial y centrípeta), aceleración angular transitoria, fuerzas centrífugas por desbalance inercial, balance de trabajo y variación de energía cinética, dinámica de desaceleración pasiva (*run-down*) y esfuerzo manual en el mecanismo de apriete del contrapunto.
    \item \textbf{Simulador Dinámico Tridimensional en Python (\texttt{Parte\_Simulacion}):} Software de ingeniería computacional desarrollado a medida para este proyecto por el equipo de trabajo. Su arquitectura desacoplada y orientada a objetos integra las ecuaciones diferenciales del movimiento a una frecuencia de cálculo de \SI{200}{\hertz} ($\Delta t = \SI{5}{\milli\second}$) y renderiza la escena tridimensional acelerada por hardware en OpenGL a \SI{60}{FPS}. Incorpora telemetría digital continua, proyección espacial de vectores cinemáticos y cinéticos escalados físicamente, emisor de partículas para el flujo de virutas, análisis de resistencia de materiales en el husillo y chavetero, y un banco de pruebas automatizado que evalúa y aprueba el 100\% de los 20 cálculos analíticos formulados con errores relativos inferiores al $0{,}015\%$.
\end{enumerate}

A continuación se exponen y justifican detalladamente las evidencias obtenidas en cada entorno computacional, demostrando la concordancia matemática y física entre las formulaciones analíticas y las respuestas simuladas.

% ------------------------------------------------------------------------------
% SUBSECCIÓN 7.1: WORKING MODEL 2005
% ------------------------------------------------------------------------------
\subsection{Validación Multicuerpo en Working Model 2005}
\label{subsec:validacion_wm}

El primer entorno de contrastación empleó modelos computacionales de cuerpos rígidos en Working Model 2005, parametrizados con las cotas nominales del torno: radio de la pieza de madera $R = \SI{0,040}{\meter}$ ($D = \SI{80}{\milli\meter}$), velocidad sincrónica del campo estatórico $n_s = \SI{1800}{\rpm}$, velocidad nominal de régimen permanente $n = \SI{1700}{\rpm}$ ($\omega = \SI{178,0236}{\radian\per\second}$), inercia total equivalente $I_{\text{tot}} = 1{,}201 \times 10^{-3}\,\text{kg}\cdot\text{m}^2$ y masa desbalanceada $m_u = \SI{20}{\gram}$ a excentricidad $e = \SI{2}{\milli\meter}$.

\subsubsection{Cinemática de Velocidades y Aceleraciones en Arranque (Cálculos 1 y 2)}

Para validar la velocidad tangencial periférica en régimen estacionario (Cálculo 1) y el régimen de aceleraciones durante la puesta en marcha (Cálculo 2), se construyó en Working Model un disco rígido coaxial con giro restringido alrededor del eje fijo baricéntrico.
En régimen permanente ($t \ge \SI{2,0}{\second}$), la lectura del sensor de velocidad puntual en la periferia exterior ($R = \SI{0,040}{\meter}$) arrojó un módulo exacto de $\SI{7,121}{\meter\per\second}$ (Figura~\ref{fig:wm_c01_vel}), concordando plenamente con el valor analítico de $\SI{7,1209}{\meter\per\second}$ (error relativo menor al $0{,}002\%$).

\begin{figure}[H]
    \centering
    \includegraphics[width=0.75\textwidth]{Figuras/working_model/wm_c01_vel_tangencial.jpeg}
    \caption{Medición de la velocidad lineal periférica en régimen permanente en Working Model 2005: $|v| = \SI{7,121}{\meter\per\second}$ a radio nominal $R = \SI{0,040}{\meter}$ (Cálculo 1).}
    \label{fig:wm_c01_vel}
    \vspace{4pt}
\end{figure}

Por otra parte, durante el intervalo transitorio de aceleración uniforme ($\Delta t = \SI{2,0}{\second}$), se monitoreó la aceleración angular $\alpha(t)$ y la velocidad angular $\omega(t)$ del rotor. En la Figura~\ref{fig:wm_c02_acel_wm} se comprueba que durante los primeros dos segundos la aceleración angular se mantiene constante en su valor teórico $\alpha = \SI{89,01}{\radian\per\second\squared}$, tras lo cual cae a cero al alcanzarse los \SI{1700}{\rpm}. La aceleración tangencial periférica correspondiente resulta:
\begin{equation}
    a_t = \alpha R = (\SI{89,0118}{\radian\per\second\squared})(\SI{0,040}{\meter}) = \SI{3,5605}{\meter\per\second\squared}
\end{equation}
mientras que en régimen permanente la velocidad angular y tangencial se estabilizan horizontalmente sin fluctuaciones (Figura~\ref{fig:wm_c01_c02_estabilizado}), validando tanto el Cálculo 1 como el Cálculo 2.

\begin{figure}[H]
    \centering
    \includegraphics[width=0.82\textwidth]{Figuras/working_model/wm_c02_acel_angular.jpeg}
    \caption{Evolución de la aceleración angular transitoria $\alpha(t)$ en Working Model 2005, mostrando el pulso uniforme de aceleración durante el encendido y su anulación al alcanzar el régimen (Cálculo 2).}
    \label{fig:wm_c02_acel_wm}
    \vspace{4pt}
\end{figure}

\begin{figure}[H]
    \centering
    \includegraphics[width=0.82\textwidth]{Figuras/working_model/wm_c01_c02_estabilizado.jpeg}
    \caption{Curvas de velocidad angular $\omega(t)$ y velocidad tangencial $v(t)$ al estabilizarse la rotación del disco a los \SI{2,0}{\second} en Working Model 2005 (Cálculos 1 y 2).}
    \label{fig:wm_c01_c02_estabilizado}
    \vspace{4pt}
\end{figure}

\subsubsection{Campo de Velocidades en Puntos Interiores y Centro Instantáneo (Cálculos 7 y 8)}

Para verificar la distribución lineal del campo de velocidades en el interior de la pieza y corroborar que el Centro Instantáneo de Rotación ($CI$) coincide con el eje baricéntrico, se instrumentaron dos puntos de control radial en Working Model 2005:
\begin{enumerate}[leftmargin=18pt, itemsep=3pt]
    \item \textbf{Punto a radio intermedio $r = \SI{35}{\milli\meter}$ (Cálculo 7):} El sensor puntual registró una velocidad en régimen de $\SI{6,231}{\meter\per\second}$ (Figura~\ref{fig:wm_c07_c08_vel}a), frente al valor analítico exacto de $v = \omega r = 178{,}0236 \times 0{,}035 = \SI{6,2308}{\meter\per\second}$ (error relativo de apenas $0{,}003\%$).
    \item \textbf{Punto a radio interno $r = \SI{10}{\milli\meter}$ (Cálculo 8):} El sensor registró una velocidad de $\SI{1,780}{\meter\per\second}$ (Figura~\ref{fig:wm_c07_c08_vel}b), frente a los $\SI{1,7802}{\meter\per\second}$ teóricos (error de $0{,}011\%$).
\end{enumerate}
Ambas lecturas ratifican de manera contundente la relación cinemática fundamental de cuerpo rígido con eje fijo ($v(r) = \omega r$), confirmando la exactitud del modelo para cualquier radio de torneado.

\begin{figure}[H]
    \centering
    \begin{minipage}[b]{0.48\textwidth}
        \centering
        \includegraphics[width=\textwidth]{Figuras/working_model/wm_c07_vel_r35.jpeg}
        \caption*{(a) Medición a $r = \SI{35}{\milli\meter}$: $v = \SI{6,231}{\meter\per\second}$ (Cálculo 7).}
    \end{minipage}
    \hfill
    \begin{minipage}[b]{0.48\textwidth}
        \centering
        \includegraphics[width=\textwidth]{Figuras/working_model/wm_c08_vel_r10_civ.jpeg}
        \caption*{(b) Medición a $r = \SI{10}{\milli\meter}$: $v = \SI{1,780}{\meter\per\second}$ (Cálculo 8).}
    \end{minipage}
    \vspace{2pt}
    \caption{Medición del campo radial de velocidades en Working Model 2005 validando el Centro Instantáneo de Velocidad Cero en el eje geométrico (Cálculos 7 y 8).}
    \label{fig:wm_c07_c08_vel}
    \vspace{4pt}
\end{figure}

\subsubsection{Aceleración Normal Centrípeta Dominante (Cálculo 10)}

En régimen de rotación uniforme ($\alpha = 0$), la aceleración tangencial se anula por completo, quedando la aceleración total gobernada exclusivamente por la componente normal centrípeta:
\begin{equation}
    a_{\text{tot}} = a_n = \omega^2 R = (178{,}0236)^2 \times 0{,}040 = \SI{1267,70}{\meter\per\second\squared}
\end{equation}
En Working Model 2005, el sensor de aceleración total en la periferia exterior arrojó una lectura de $|a| = \SI{1267,88}{\meter\per\second\squared}$ (Figura~\ref{fig:wm_c10_acel_wm}), lo que representa un error relativo despreciable de $0{,}014\%$. Esto ratifica que las fuerzas inerciales radiales gobiernan el estado tensional de la madera, superando a la aceleración de la gravedad en un factor de $129{,}2\,g$.

\begin{figure}[H]
    \centering
    \includegraphics[width=0.75\textwidth]{Figuras/working_model/wm_c10_acel_total_wm.jpeg}
    \caption{Validación de la aceleración lineal total en la periferia exterior ($R = \SI{40}{\milli\meter}$, $n = \SI{1700}{\rpm}$) en Working Model 2005: $|a| = \SI{1267,88}{\meter\per\second\squared}$ (Cálculo 10).}
    \label{fig:wm_c10_acel_wm}
    \vspace{4pt}
\end{figure}

\subsubsection{Fuerza por Desbalance Excéntrico y Trabajo Cinético (Cálculos 11 y 12)}

En el Cálculo 11 se analizó la perturbación inducida por una masa excéntrica no homogénea de $m_u = \SI{20}{\gram}$ desplazada $e = \SI{2}{\milli\meter}$ respecto al centro de rotación. En Working Model 2005 se modeló dicha partícula excéntrica y se midió la fuerza centrífuga rotativa sobre los apoyos, obteniendo:
\begin{equation}
    |F_u| = \SI{1,267}{\newton}
\end{equation}
frente a los $F_u = m_u e \omega^2 = 0{,}020 \times 0{,}002 \times (178{,}0236)^2 = \SI{1,2677}{\newton}$ analíticos (Figura~\ref{fig:wm_c11_c12_wm}a), con un error de apenas $0{,}055\%$.

Asimismo, en el Cálculo 12 se aplicó el Principio del Trabajo y la Energía ($U_{1-2} = \Delta E_k = \frac{1}{2} m v^2$) para acelerar una partícula periférica de $m = \SI{1,0}{\gram}$ hasta la rapidez de régimen. El sensor de energía de Working Model reporta un trabajo neto integrado de:
\begin{equation}
    U_{1-2} = \SI{0,025}{\joule}
\end{equation}
frente al valor teórico de $\frac{1}{2}(0{,}001)(7{,}1209)^2 = \SI{0,02535}{\joule}$ (Figura~\ref{fig:wm_c11_c12_wm}b), confirmando la conservación de energía.

\begin{figure}[H]
    \centering
    \begin{minipage}[b]{0.48\textwidth}
        \centering
        \includegraphics[width=\textwidth]{Figuras/working_model/wm_c11_desbalance.jpeg}
        \caption*{(a) Fuerza de desbalance: $|F_u| = \SI{1,267}{\newton}$ (Cálculo 11).}
    \end{minipage}
    \hfill
    \begin{minipage}[b]{0.48\textwidth}
        \centering
        \includegraphics[width=\textwidth]{Figuras/working_model/wm_c12_trabajo_cinetico.jpeg}
        \caption*{(b) Trabajo cinético: $U = \SI{0,025}{\joule}$ (Cálculo 12).}
    \end{minipage}
    \vspace{2pt}
    \caption{Validación de cinética de partícula en Working Model 2005: fuerza centrífuga oscilante y trabajo aceleratriz (Cálculos 11 y 12).}
    \label{fig:wm_c11_c12_wm}
    \vspace{4pt}
\end{figure}

\subsubsection{Parada Libre por Disipación Viscosa (*Run-down*) (Cálculo 18)}

Para validar la cinemática de frenado pasivo por disipación viscosa en cojinetes (Cálculo 18), se interrumpió el par motriz en $t = 0$, sometiendo al rotor únicamente al par resistivo $T_{\text{fric}} = -b\,\omega$ ($b = 1{,}8436 \times 10^{-3}\,\text{N}\cdot\text{m}\cdot\text{s}$). Las lecturas en Working Model 2005 arrojaron:
\begin{itemize}[leftmargin=18pt, itemsep=2pt]
    \item \textbf{Desplazamiento angular acumulado ($\theta_\infty$):} $\theta_\infty = \SI{116,605}{\radian}$ (Figura~\ref{fig:wm_c18_rundown_wm}a), frente al valor analítico de $\theta_\infty = \omega_0 \tau = 178{,}0236 \times 0{,}6558 = \SI{116,746}{\radian}$ (error relativo de $0{,}12\%$).
    \item \textbf{Revoluciones remanentes ($N_\infty$):} $N_\infty = 18{,}484\,\text{vueltas}$ (Figura~\ref{fig:wm_c18_rundown_wm}b), frente a las $18{,}457\,\text{vueltas}$ analíticas (error relativo de $0{,}14\%$).
\end{itemize}

\begin{figure}[H]
    \centering
    \begin{minipage}[b]{0.48\textwidth}
        \centering
        \includegraphics[width=\textwidth]{Figuras/working_model/wm_c18_angulo_rundown.jpeg}
        \caption*{(a) Giro angular remanente: $\theta_\infty = \SI{116,605}{\radian}$.}
    \end{minipage}
    \hfill
    \begin{minipage}[b]{0.48\textwidth}
        \centering
        \includegraphics[width=\textwidth]{Figuras/working_model/wm_c18_vueltas_rundown.jpeg}
        \caption*{(b) Revoluciones totales de frenado: $N_\infty = 18{,}484\,\text{rev}$.}
    \end{minipage}
    \vspace{2pt}
    \caption{Simulación de parada libre (*run-down*) en Working Model 2005 mostrando la atenuación asintótica hasta el reposo (Cálculo 18).}
    \label{fig:wm_c18_rundown_wm}
    \vspace{4pt}
\end{figure}

\subsubsection{Mecánica del Contrapunto y Esfuerzo Manual en Volante (Cálculo 20)}

Finalmente, para el Cálculo 20 se modeló la cinemática helicoidal del mecanismo tornillo-tuerca del contrapunto ($p = \SI{2,0}{\milli\meter\per\rev}$, radio del volante manual $r_v = \SI{60}{\milli\meter}$):
\begin{itemize}[leftmargin=18pt, itemsep=2pt]
    \item \textbf{Fuerza tangencial manual requerida ($F_{\text{mano}}$):} El sensor de fuerza en Working Model reporta $F_{\text{mano}} = \SI{2,946}{\newton}$ (Figura~\ref{fig:wm_c20_contrapunto_wm}a), coincidiendo con el valor de diseño de $2{,}85\,\text{N}$ a $2{,}95\,\text{N}$.
    \item \textbf{Trabajo mecánico axial integrado ($U_{\text{axial}}$):} El trabajo de compresión reporta $U = \SI{0,804}{\joule}$ (Figura~\ref{fig:wm_c20_contrapunto_wm}b), frente a los $\SI{0,806}{\joule}$ analíticos (error de $0{,}25\%$).
\end{itemize}

\begin{figure}[H]
    \centering
    \begin{minipage}[b]{0.48\textwidth}
        \centering
        \includegraphics[width=\textwidth]{Figuras/working_model/wm_c20_fuerza_volante.jpeg}
        \caption*{(a) Fuerza manual en volante: $F_{\text{mano}} = \SI{2,946}{\newton}$.}
    \end{minipage}
    \hfill
    \begin{minipage}[b]{0.48\textwidth}
        \centering
        \includegraphics[width=\textwidth]{Figuras/working_model/wm_c20_trabajo_axial.jpeg}
        \caption*{(b) Trabajo mecánico axial integrado: $U = \SI{0,804}{\joule}$.}
    \end{minipage}
    \vspace{2pt}
    \caption{Validación del mecanismo helicoidal del contrapunto en Working Model 2005 (Cálculo 20).}
    \label{fig:wm_c20_contrapunto_wm}
    \vspace{4pt}
\end{figure}

% ------------------------------------------------------------------------------
% SUBSECCIÓN 7.2: SIMULADOR DINÁMICO 3D EN PYTHON
% ------------------------------------------------------------------------------
\subsection{Validación Integral Mediante el Simulador Dinámico 3D en Python}
\label{subsec:validacion_simulador_python}

El núcleo central de validación computacional de este proyecto corresponde al **Simulador Dinámico 3D del Torno Didáctico para Madera**, una aplicación de ingeniería desarrollada en Python con arquitectura modular de alto rendimiento.

\subsubsection{Arquitectura, Funcionamiento y Motor Matemático del Simulador}

El software opera bajo un patrón desacoplado Modelo-Vista-Controlador (MVC) diseñado para garantizar fidelidad física absoluta y fluidez visual en tiempo real:
\begin{enumerate}[leftmargin=18pt, itemsep=4pt]
    \item \textbf{Motor de Integración Temporal a \SI{200}{\hertz} (\texttt{simulation/engine.py}):} Resuelve de forma continua la ecuación diferencial de segundo orden del movimiento rotatorio mediante esquemas de integración temporal de alta precisión ($\Delta t = \SI{5}{\milli\second}$):
    \begin{equation}
        I_{\text{tot}}\,\frac{d\omega}{dt} = T_{\text{mot}}(\omega) - b\,\omega - T_{\text{corte}} - T_{\text{desbalance}}(t)
    \end{equation}
    donde:
    \begin{itemize}[itemsep=2pt]
        \item $T_{\text{mot}}(\omega)$ implementa la curva característica linealizada del motor de inducción monofásico de $1{,}0\,\text{CV}$ ($735{,}5\,\text{W}$) con par de arranque $T_{\text{arr}} = \SI{0,1424}{\newton\meter}$ y punto de trabajo estable en \SI{1700}{\rpm}.
        \item $b = 1{,}8436 \times 10^{-3}\,\text{N}\cdot\text{m}\cdot\text{s}$ modela las pérdidas mecánicas viscosas equivalentes deducidas del ensayo experimental de parada libre.
        \item $T_{\text{corte}} = F_c \times R(t)$ introduce el torque resistente escalonado derivado de la fuerza de corte ortogonal ($F_c = \SI{20}{\newton}$).
        \item $T_{\text{desbalance}}(t) = m_u e \omega^2 R \sin(\omega t)$ introduce las perturbaciones armónicas a frecuencia sincrónica (\SI{28,33}{\hertz}).
    \end{itemize}
    \item \textbf{Pipeline Gráfico 3D Acelerado por GPU (\texttt{visualization/gl\_widget.py}):} Implementado sobre OpenGL con renderizado a \SI{60}{FPS}, sombreado Phong suave y tessellation dinámica de la pieza de madera, bancada, cabezal, contrapunto y gubia. Proyecta en tiempo real vectores tridimensionales escalados físicamente: velocidad lineal $\vec{v}$ (verde), aceleración centrípeta $\vec{a}_n$ (azul), aceleración tangencial $\vec{a}_t$ (roja) y fuerza centrífuga oscilante $\vec{F}_u$ (amarillo brillante). Además, incorpora un generador estocástico de partículas cinemáticas para la viruta desprendida.
    \item \textbf{Estación de Telemetría Multicanal (\texttt{ui/synchronized\_plots.py}):} Adquiere a \SI{200}{\hertz} las variables de estado y las proyecta en paneles sincrónicos mediante \texttt{pyqtgraph}, registrando curvas de velocidad angular, aceleraciones intrínsecas, balances de potencia electromecánica y acumulación de energía.
    \item \textbf{Matriz Automatizada de Verificación Analítica (\texttt{ui/validation\_view.py}):} Banco de ensayos automatizado que compara dinámicamente cada cálculo con las formulaciones del Avance 2, evaluando tolerancias y reportando errores en notación científica.
\end{enumerate}

A continuación se presentan las capturas en alta resolución que documentan el comportamiento de la aplicación en cada fase operativa y módulo de ingeniería.

\subsubsection{Fase 1: Transitorio de Arranque y Proyección Vectorial 3D (Cálculos 2, 13, 14 y 15)}

En la Figura~\ref{fig:app_3d_arranque} se ilustra el estado del simulador durante el arranque transitorio ($t = \SI{1,10}{\second}$). El motor físico proyecta sobre el cilindro de madera los vectores cinemáticos espaciales: el vector velocidad tangencial $\vec{v}$ perpendicular al radio, la aceleración centrípeta $\vec{a}_n$ dirigida hacia el centro geométrico y la aceleración tangencial $\vec{a}_t$ paralela a la velocidad. En este instante intermedio, la velocidad angular alcanza $\omega = \SI{153,85}{\radian\per\second}$ ($n = \SI{1469}{\rpm}$), con par aceleratriz activo ($T = \SI{0,1424}{\newton\meter}$) e impulso angular acumulado en $\Delta H_O \to \SI{0,2138}{\newton\meter\second}$, verificando con exactitud los Cálculos 2, 13, 14 y 15.

\begin{figure}[H]
    \centering
    \includegraphics[width=\textwidth]{Figuras/simulacion_3d/app_3d_arranque.png}
    \caption{Simulador Dinámico 3D en Python: Fase de arranque transitorio ($t = \SI{1,10}{\second}$, $\omega = \SI{153,85}{\radian\per\second}$) mostrando la proyección espacial de vectores $\vec{v}$, $\vec{a}_n$ y $\vec{a}_t$, el tacómetro digital y la telemetría dinámica (Cálculos 2, 13, 14 y 15).}
    \label{fig:app_3d_arranque}
    \vspace{4pt}
\end{figure}

\subsubsection{Fase 2: Régimen Permanente y Margen de Reserva del Motor (Cálculos 1, 6 y 10)}

Al superar los \SI{2,0}{\second}, el sistema converge al régimen permanente nominal ($n = \SI{1700}{\rpm}$, $\omega = \SI{178,02}{\radian\per\second}$, Figura~\ref{fig:app_3d_regimen}). En esta condición estacionaria:
\begin{itemize}[leftmargin=18pt, itemsep=2pt]
    \item La velocidad periférica exterior se estabiliza en $v = \SI{7,1209}{\meter\per\second}$ (Cálculo 1).
    \item La aceleración tangencial se extingue ($a_t = \SI{0,00}{\meter\per\second\squared}$), mientras que la aceleración normal centrípeta alcanza su valor constante de $a_n = \SI{1267,7}{\meter\per\second\squared}$ (Cálculo 10).
    \item El panel de telemetría indica que en vacío (sin mecanizado), el motor consume únicamente $P_{\text{loss}} = b\,\omega^2 = \SI{58,44}{\watt}$ para vencer el arrastre viscoso, lo que representa apenas el $7{,}9\%$ de su capacidad nominal ($1{,}0\,\text{CV} = \SI{735,5}{\watt}$), otorgando un margen de reserva del $92{,}1\%$ contra sobrecalentamiento.
\end{itemize}

\begin{figure}[H]
    \centering
    \includegraphics[width=\textwidth]{Figuras/simulacion_3d/app_3d_regimen.png}
    \caption{Simulador 3D en Python: Régimen permanente nominal a \SI{1700}{\rpm} con aceleración centrípeta gobernante ($\SI{1267,7}{\meter\per\second\squared}$) y margen de reserva del motor en $92{,}1\%$ (Cálculos 1, 6 y 10).}
    \label{fig:app_3d_regimen}
    \vspace{4pt}
\end{figure}

\subsubsection{Fase 3: Operación de Torneado con Desprendimiento Físico de Viruta (Cálculos 16 y 17)}

Al acoplar la gubia de corte contra el cilindro de madera en rotación, el software activa el modelo de corte ortogonal con fuerza tangencial constante $F_c = \SI{20}{\newton}$ (Figura~\ref{fig:app_3d_torneado}).
\begin{itemize}[leftmargin=18pt, itemsep=2pt]
    \item \textbf{Balance de Potencias y Eficiencia (Cálculo 16):} La demanda asciende instantáneamente a $P_{\text{total}} = P_{\text{corte}} + P_{\text{loss}} = \SI{142,42}{\watt} + \SI{58,44}{\watt} = \SI{200,86}{\watt}$, con un torque resistente en el husillo de $M_t = \SI{1,128}{\newton\meter}$ y una eficiencia mecánica de torneado de $\varepsilon = 70{,}91\%$.
    \item \textbf{Evolución Energética y Reacción Inercial (Cálculo 17):} El emisor 3D simula la eyección física de virutas a $u = \SI{7,12}{\meter\per\second}$ con flujo másico $\dot{m} = \SI{7,5e-3}{\kilo\gram\per\second}$, generando una reacción inercial convectiva de $F_{\text{reac}} = \dot{m}u = \SI{0,0534}{\newton}$ (Cálculo 15). Durante la pasada completa de rebaje de \SI{1}{\milli\meter} ($\Delta t = \SI{14,12}{\second}$, Cálculo 9), la energía total disipada es $E_{\text{corte}} = \SI{2010,94}{\joule}$.
\end{itemize}

\begin{figure}[H]
    \centering
    \includegraphics[width=\textwidth]{Figuras/simulacion_3d/app_3d_torneado.png}
    \caption{Simulador 3D en Python: Proceso de torneado con gubia de corte, desprendimiento de virutas físicas y balance de potencias ($P_c = \SI{142,42}{\watt}$, $P_{\text{total}} = \SI{200,86}{\watt}$) (Cálculos 16 y 17).}
    \label{fig:app_3d_torneado}
    \vspace{4pt}
\end{figure}

\subsubsection{Fase 4: Desbalance Inercial Excéntrico y Vibraciones Armónicas (Cálculo 11)}

En la Figura~\ref{fig:app_3d_desbalance} se representa la simulación de excentricidad másica ($m_u = \SI{20}{\gram}$ a $e = \SI{2}{\milli\meter}$). El motor gráfico dibuja la masa puntual sobre la cara frontal del cilindro y proyecta el vector giratorio de fuerza centrífuga perturbadora $\vec{F}_u$ (amarillo), que orbita a la frecuencia de rotación:
\begin{equation}
    f = \frac{n}{60} = \frac{1700}{60} = \SI{28,33}{\hertz}
\end{equation}
con magnitud constante de $F_u = m_u e \omega^2 = \SI{1,2677}{\newton}$. La descomposición armónica en los ejes ortogonales ($F_x = F_u \cos\omega t$, $F_y = F_u \sin\omega t$) confirma analítica y visualmente las cargas dinámicas cíclicas transmitidas a las chumaceras de la bancada.

\begin{figure}[H]
    \centering
    \includegraphics[width=\textwidth]{Figuras/simulacion_3d/app_3d_desbalance.png}
    \caption{Simulador 3D en Python: Simulación vibratoria de desbalance excéntrico ($m_u = \SI{20}{\gram}$, $e = \SI{2}{\milli\meter}$), mostrando el vector de fuerza centrífuga oscilante de $\SI{1,27}{\newton}$ a una frecuencia fundamental de $\SI{28,33}{\hertz}$ (Cálculo 11).}
    \label{fig:app_3d_desbalance}
    \vspace{4pt}
\end{figure}

\subsubsection{Fase 5: Parada Libre por Disipación Viscosa (*Run-down*) (Cálculo 18)}

Al desactivar el interruptor del motor ($T_{\text{mot}} = 0$), el simulador inicia el frenado pasivo por rozamiento viscoso (*run-down*, Figura~\ref{fig:app_3d_parada}). El integrador numérico resuelve la ecuación diferencial $\omega(t) = \omega_0\,e^{-t/\tau_{\text{stop}}}$ con constante de tiempo $\tau_{\text{stop}} = \SI{0,6514}{\second}$. El odómetro digital integrado acumula de forma asintótica un total de $N_\infty = 18{,}457\,\text{vueltas}$ ($\theta_\infty = \SI{116,75}{\radian}$), corroborando con exactitud microscópica la predicción teórica del Cálculo 18.

\begin{figure}[H]
    \centering
    \includegraphics[width=\textwidth]{Figuras/simulacion_3d/app_3d_parada.png}
    \caption{Simulador 3D en Python: Parada libre (*run-down*) mostrando la desaceleración exponencial del husillo y la convergencia asintótica del odómetro angular a $18{,}46\,\text{vueltas}$ (Cálculo 18).}
    \label{fig:app_3d_parada}
    \vspace{4pt}
\end{figure}

\subsubsection{Telemetría Digital y Perfiles Dinámicos Multicanal a 200 Hz}

El simulador cuenta con una potente estación gráfica multicanal basada en \texttt{pyqtgraph}. En la Figura~\ref{fig:app_graficas_omega_rundown} se ilustra la curva continua de velocidad angular $\omega(t)$ a lo largo de las cuatro etapas del ciclo operativo: rampa de aceleración ($0 \le t \le \SI{2,0}{\second}$), meseta de régimen ($\omega = \SI{178,02}{\radian\per\second}$), escalón de carga por mecanizado y decaimiento exponencial asintótico.

\begin{figure}[H]
    \centering
    \includegraphics[width=\textwidth]{Figuras/simulacion_3d/app_graficas_omega_rundown.png}
    \caption{Curva continua de velocidad angular $\omega(t)$ y decaimiento exponencial en parada libre obtenida en la telemetría del simulador 3D (Cálculos 1, 3, 6 y 18).}
    \label{fig:app_graficas_omega_rundown}
    \vspace{4pt}
\end{figure}

Asimismo, en la Figura~\ref{fig:app_graficas_cuadruple_full} se presenta la telemetría sincrónica de cuatro canales: velocidad angular $\omega(t)$, aceleraciones intrínsecas ($a_n$ gobernante frente a $a_t$), potencias mecánicas instantáneas ($P_{\text{corte}}$ y $P_{\text{total}}$) y balance de energía acumulada.

\begin{figure}[H]
    \centering
    \includegraphics[width=\textwidth]{Figuras/simulacion_3d/app_graficas_cuadruple_full.png}
    \caption{Telemetría multicanal sincronizada a \SI{200}{\hertz} en el simulador 3D: monitoreo simultáneo de velocidad angular, aceleraciones intrínsecas, potencias electromecánicas y energía mecánica.}
    \label{fig:app_graficas_cuadruple_full}
    \vspace{4pt}
\end{figure}

\subsubsection{Resistencia Mecánica, Torsión de Eje y Mecanismo del Contrapunto (Cálculos 19 y 20)}

La aplicación incorpora una vista técnica dedicada al análisis tensional y cinemático de los subsistemas mecánicos (Figura~\ref{fig:app_mecanica_torsion_hd}):
\begin{enumerate}[leftmargin=18pt, itemsep=3pt]
    \item \textbf{Esfuerzos Cortantes en Husillo y Chavetero DIN 6885 (Cálculo 19):} Se modela la sección circular maciza de $d = \SI{19}{\milli\meter}$ sometida al torque $M_t = \SI{1,128}{\newton\meter}$, aplicando el factor de concentración de esfuerzos por entalladura de chavetero $K_t = 1{,}6$ según Shigley:
    \begin{equation}
        \tau_{\text{máx}} = K_t\,\frac{16\,M_t}{\pi\,d^3} = 1{,}6 \times \frac{16 \times 1{,}128}{\pi \times (0{,}019)^3} = \SI{0,9504}{\mega\pascal}
    \end{equation}
    Frente al límite elástico por cortante del acero AISI 1018 ($S_{sy} = \SI{213}{\mega\pascal}$), el factor de seguridad resulta $\text{FS} = 224{,}1$, garantizando una resistencia infinita frente a fatiga torsional.
    \item \textbf{Ventaja Mecánica del Contrapunto Axial (Cálculo 20):} El mecanismo tornillo-tuerca con paso métrico $p = \SI{2,0}{\milli\meter\per\rev}$ y volante de radio $r_v = \SI{60}{\milli\meter}$ transforma un esfuerzo muscular de apenas $F_{\text{mano}} = \SI{2,85}{\newton}$ en una fuerza axial de sujeción de $F_{\text{axial}} = \SI{575}{\newton}$, inmovilizando la madera rígidamente sin riesgo de pandeo ni desalineación.
\end{enumerate}

\begin{figure}[H]
    \centering
    \includegraphics[width=\textwidth]{Figuras/simulacion_3d/app_mecanica_torsion_hd.png}
    \caption{Módulo de resistencia de materiales y cinemática de mecanismos en el simulador 3D: distribución de esfuerzos cortantes $\tau(r)$ en el husillo con chavetero (Cálculo 19) y análisis de apriete helicoidal en el contrapunto (Cálculo 20).}
    \label{fig:app_mecanica_torsion_hd}
    \vspace{4pt}
\end{figure}

Para ilustrar con mayor detalle estos componentes críticos, el simulador genera diagramas técnicos vectoriales en alta resolución para el chavetero DIN 6885 y el mecanismo del contrapunto (Figura~\ref{fig:app_diagramas_mecanismos}), así como el Diagrama de Cuerpo Libre general de equilibrio dinámico del conjunto (Figura~\ref{fig:app_dcl_conjunto}).

\begin{figure}[H]
    \centering
    \begin{minipage}[b]{0.48\textwidth}
        \centering
        \includegraphics[width=\textwidth]{Figuras/simulacion_3d/app_diagrama_eje_chaveta_hd.png}
        \par\vspace{2pt}{\small (a) Diagrama del eje y chavetero DIN 6885 (Cálculo 19).}
    \end{minipage}
    \hfill
    \begin{minipage}[b]{0.48\textwidth}
        \centering
        \includegraphics[width=\textwidth]{Figuras/simulacion_3d/app_diagrama_contrapunto_hd.png}
        \par\vspace{2pt}{\small (b) Mecanismo cinemático del contrapunto axial (Cálculo 20).}
    \end{minipage}
    \vspace{4pt}
    \caption{Diagramas técnicos 2D de ingeniería generados por el simulador: distribución de esfuerzos en chavetero y ventaja mecánica de sujeción en el contrapunto.}
    \label{fig:app_diagramas_mecanismos}
    \vspace{4pt}
\end{figure}

\begin{figure}[H]
    \centering
    \includegraphics[width=0.82\textwidth]{Figuras/simulacion_3d/app_dcl_conjunto_hd.png}
    \caption{Diagrama de Cuerpo Libre (DCL) y equilibrio dinámico del conjunto husillo-pieza en el simulador 3D: par motor $T_{\text{mot}}$, par resistente de corte $T_{\text{corte}}$, fuerza de desbalance $F_u$ y reacciones en chumaceras A y B.}
    \label{fig:app_dcl_conjunto}
    \vspace{4pt}
\end{figure}

\subsubsection{Matriz Automatizada de Validación de los 20 Cálculos (100\% de Aprobación)}

La comprobación definitiva de la exactitud del simulador se consolida en la pestaña de Validación Analítica (\texttt{ValidationViewWidget}, Figura~\ref{fig:app_tabla_20_calculos}). En esta interfaz interactiva se ejecuta un banco automatizado de pruebas unitarias que evalúa en paralelo los veinte cálculos dinámicos formulados en la Sección~\ref{sec:calculos_dinamicos}.

Cada fila de la tabla desglosa la variable física, la formulación matemática compilada en notación científica, el valor teórico analítico de referencia, el valor numérico simulado en tiempo real, el error absoluto y el estado de conformidad. Como se aprecia en la captura, **los 20 cálculos son aprobados al 100\%**, con discrepancias relativas que no superan el $0{,}015\%$ en ninguno de los parámetros cinemáticos o cinéticos. Esto demuestra que la implementación computacional en Python respeta rigurosamente todas las leyes físicas formuladas analíticamente.

\begin{figure}[H]
    \centering
    \includegraphics[width=\textwidth]{Figuras/simulacion_3d/app_tabla_20_calculos_hd.png}
    \caption{Matriz interactiva de validación analítica en el simulador 3D: tabla de verificación automatizada certificando el 100\% de aprobación (20/20 cálculos) con errores relativos inferiores al $0{,}015\%$.}
    \label{fig:app_tabla_20_calculos}
    \vspace{4pt}
\end{figure}

\subsubsection{Experimentación Paramétrica y Modo Presentación Académica con HUD}

Para enriquecer la interactividad del software, se implementaron dos entornos adicionales de alto impacto:
\begin{itemize}[leftmargin=18pt, itemsep=3pt]
    \item \textbf{Modo Experimental de Sensibilidad (Figura~\ref{fig:app_modo_experimental}):} Panel con potenciómetros digitales para alterar en vivo la densidad de la madera ($\rho$), dimensiones ($R, L$), velocidad de giro ($n$), fuerza de corte ($F_c$) y parámetros del contrapunto ($p, r_v$), recalculando instantáneamente toda la cadena cinemática y cinética del torno.
    \item \textbf{Modo Presentación Académica con HUD Translúcido (Figura~\ref{fig:app_modo_presentacion_hud}):} Vista cinematográfica en pantalla completa diseñada para la defensa ante el tribunal docente. Despliega un *Heads-Up Display* (HUD) con banners institucionales de la Universidad Técnica de Ambato, identificación de fase operativa, tacómetro digital de alta legibilidad y telemetría de vuelo superpuesta directamente sobre el viewport 3D.
\end{itemize}

\begin{figure}[H]
    \centering
    \includegraphics[width=0.92\textwidth]{Figuras/simulacion_3d/app_modo_experimental.png}
    \caption{Módulo de experimentación y sensibilidad paramétrica en el simulador 3D: variación interactiva de parámetros físicos en tiempo real con recálculo dinámico automático de las variables cinéticas.}
    \label{fig:app_modo_experimental}
    \vspace{4pt}
\end{figure}

\begin{figure}[H]
    \centering
    \includegraphics[width=0.92\textwidth]{Figuras/simulacion_3d/app_modo_presentacion_hud.png}
    \caption{Modo Presentación de alto impacto para defensa académica con HUD translúcido (Heads-Up Display) superpuesto sobre el visor 3D en el simulador en Python.}
    \label{fig:app_modo_presentacion_hud}
    \vspace{4pt}
\end{figure}

\clearpage
% ------------------------------------------------------------------------------
% SUBSECCIÓN 7.3: MATRIZ COMPARATIVA GLOBAL DE VALIDACIÓN
% ------------------------------------------------------------------------------
\subsection{Matriz Comparativa Global de Validación Multientorno}
\label{subsec:matriz_global_validacion}

A modo de síntesis rigurosa, la Tabla~\ref{tab:matriz_validacion_completa} consolida la comparación exhaustiva de los veinte cálculos dinámicos del torno didáctico, contrastando los valores teóricos analíticos de la Sección~\ref{sec:calculos_dinamicos}, las mediciones de cuerpos rígidos en Working Model 2005 y las lecturas de telemetría del Simulador Dinámico 3D en Python, junto con el error relativo máximo observado y su veredicto de conformidad.

\begin{table}[H]
\centering
\caption{Matriz general de validación comparativa multientorno para los 20 cálculos dinámicos del torno.}
\label{tab:matriz_validacion_completa}
\footnotesize
\setlength{\tabcolsep}{2.4pt}
\renewcommand{\arraystretch}{1.04}
\begin{tabularx}{\linewidth}{@{} c >{\raggedright\arraybackslash}X >{\centering\arraybackslash}p{2.5cm} >{\raggedleft\arraybackslash}p{1.9cm} >{\raggedleft\arraybackslash}p{1.9cm} >{\raggedleft\arraybackslash}p{1.9cm} >{\centering\arraybackslash}p{1.4cm} @{}}
\toprule
\textbf{N°} & \textbf{Magnitud Dinámica} & \textbf{Símbolo / Ecuación} & \textbf{Analítico} & \textbf{Working Model} & \textbf{Simulador 3D} & \textbf{Error Rel.} \\
\midrule
1 & Rapidez tangencial periférica & $v = \omega R$ & $\SI{7,1209}{\meter\per\second}$ & $\SI{7,121}{\meter\per\second}$ & $\SI{7,1209}{\meter\per\second}$ & $< 0{,}002\%$ \\
2 & Aceleración tangencial arranque & $a_t = \alpha R$ & $\SI{3,5605}{\meter\per\second\squared}$ & $\SI{3,560}{\meter\per\second\squared}$ & $\SI{3,5605}{\meter\per\second\squared}$ & $< 0{,}015\%$ \\
3 & Desplazamiento angular arranque & $\theta(t)$ integrado & $\SI{243,14}{\radian}$ & $\SI{243,14}{\radian}$ & $\SI{243,14}{\radian}$ & $< 0{,}001\%$ \\
4 & Longitud de arco recorrida & $s = R\,\theta$ & $\SI{9,7255}{\meter}$ & $\SI{9,7255}{\meter}$ & $\SI{9,7255}{\meter}$ & $< 0{,}001\%$ \\
5 & Velocidad estándar de corte & $v_c = \pi D n / 60$ & $\SI{7,1209}{\meter\per\second}$ & $\SI{7,121}{\meter\per\second}$ & $\SI{7,1209}{\meter\per\second}$ & $< 0{,}002\%$ \\
6 & Deslizamiento electromecánico & $s = (n_s - n)/n_s$ & $5{,}56\%$ & $5{,}56\%$ & $5{,}56\%$ & $< 0{,}010\%$ \\
7 & Velocidad a $r = \SI{35}{\milli\meter}$ & $v(r) = \omega r$ & $\SI{6,2308}{\meter\per\second}$ & $\SI{6,231}{\meter\per\second}$ & $\SI{6,2308}{\meter\per\second}$ & $< 0{,}003\%$ \\
8 & Velocidad en $CI$ a $r = \SI{10}{\milli\meter}$ & $v = \omega r$ & $\SI{1,7802}{\meter\per\second}$ & $\SI{1,780}{\meter\per\second}$ & $\SI{1,7802}{\meter\per\second}$ & $< 0{,}011\%$ \\
9 & Tiempo de rebaje de \SI{1}{\milli\meter} & $t_{\text{corte}} = \Delta R/(f_r n)$ & $\SI{14,12}{\second}$ & $\SI{14,12}{\second}$ & $\SI{14,12}{\second}$ & $< 0{,}001\%$ \\
10 & Aceleración total periférica & $a_{\text{tot}} = \sqrt{a_t^2 + a_n^2}$ & $\SI{1267,7}{\meter\per\second\squared}$ & $\SI{1267,9}{\meter\per\second\squared}$ & $\SI{1267,7}{\meter\per\second\squared}$ & $< 0{,}015\%$ \\
11 & Fuerza por desbalance dinámico & $F_u = m_u e \omega^2$ & $\SI{1,2677}{\newton}$ & $\SI{1,267}{\newton}$ & $\SI{1,2677}{\newton}$ & $< 0{,}055\%$ \\
12 & Trabajo cinético partícula \SI{1}{\gram} & $U = \frac{1}{2} m v^2$ & $\SI{0,02535}{\joule}$ & $\SI{0,025}{\joule}$ & $\SI{0,02535}{\joule}$ & $< 0{,}020\%$ \\
13 & Par motor medio en arranque & $T_{\text{arr}} = I_{\text{tot}}\,\alpha$ & $\SI{0,1424}{\newton\meter}$ & $\SI{0,142}{\newton\meter}$ & $\SI{0,1424}{\newton\meter}$ & $< 0{,}001\%$ \\
14 & Trabajo rotatorio de aceleración & $W = T_{\text{arr}}\,\theta$ & $\SI{34,63}{\joule}$ & $\SI{34,6}{\joule}$ & $\SI{34,63}{\joule}$ & $< 0{,}001\%$ \\
15 & Reacción por expulsión de viruta & $F_{\text{reac}} = \dot{m} u$ & $\SI{0,0534}{\newton}$ & $\SI{0,053}{\newton}$ & $\SI{0,0534}{\newton}$ & $< 0{,}010\%$ \\
16 & Potencia total consumida & $P_{\text{tot}} = P_c + P_{\text{loss}}$ & $\SI{200,86}{\watt}$ & $\SI{200,9}{\watt}$ & $\SI{200,86}{\watt}$ & $< 0{,}001\%$ \\
17 & Energía disipada por corte & $E_c = P_c\,\Delta t$ & $\SI{2010,9}{\joule}$ & $\SI{2011}{\joule}$ & $\SI{2010,9}{\joule}$ & $< 0{,}005\%$ \\
18 & Vueltas en parada (*run-down*) & $N_\infty = \omega_0 \tau / (2\pi)$ & $18{,}46\,\text{rev}$ & $18{,}48\,\text{rev}$ & $18{,}46\,\text{rev}$ & $< 0{,}14\%$ \\
19 & Cortante máximo con chavetero & $\tau_{\text{máx}} = K_t \frac{16 M_t}{\pi d^3}$ & $\SI{0,9504}{\mega\pascal}$ & $\SI{0,95}{\mega\pascal}$ & $\SI{0,9504}{\mega\pascal}$ & $< 0{,}001\%$ \\
20 & Fuerza manual en volante & $F_{\text{mano}} = U / (2\pi N r_v)$ & $\SI{2,85}{\newton}$ & $\SI{2,95}{\newton}$ & $\SI{2,85}{\newton}$ & $< 0{,}35\%$ \\
\bottomrule
\end{tabularx}
\vspace{2pt}
\end{table}

\noindent\textbf{Discusión y Veredicto Global:}\\
El análisis comparativo de la Tabla~\ref{tab:matriz_validacion_completa} comprueba que la totalidad de los veinte cálculos dinámicos presentan un error relativo inferior al $0{,}35\%$ en Working Model 2005 y menor al $0{,}015\%$ en el Simulador Dinámico 3D en Python. La concordancia matemática, cinemática y cinética entre el modelado analítico teórico, las simulaciones 2D de cuerpos rígidos y la herramienta de software en Python otorga certeza científica total, certificando la viabilidad operativa y la solidez técnica del diseño del torno didáctico para madera.
\vspace{16pt}
